\section{Optionaler Anhang: Reproduzierbarkeit}

Die Ausführung erfolgt mit Docker Compose. Der nachgewiesene Matrix- und
Replaypfad verwendet unmaskierte Bilder für COLMAP, hält das originale Modell
für GCP/SfM vor und führt STS sowie SuGaR in der idealen PINHOLE-Bilddomäne
aus. Bei Verzeichnungsmodellen werden die SAM-Masken dort mit derselben
Kameramodellabbildung in die ideale Domäne gewarpt. Der historische Inline-
Vollauf erzeugt die ideale Bildszene, besitzt aber noch keinen eigenen Aufruf
des Maskenwarp-Skripts.

Die wichtigsten konstanten Parameter der PA-Versuche sind:

\begin{itemize}
  \item Prompt \texttt{pipe};
  \item 4096 SIFT-Merkmale, Sequential-Overlap 15, Guided Matching aus;
  \item STS 7000 Gesamtiterationen, Stage-2-Fenster 5000;
  \item Objektmasken für die Bildmetriken und feste Eval-Listen;
  \item 200000 Mesh-Vertices, 5000000 Oberflächenstichproben und Seed 42;
  \item objektmaskierte PSNR-, SSIM- und LPIPS-Auswertung.
    \item Produktionsroute \texttt{original\_gs}; SuGaR-Coarse nur als
      explizite Vergleichsroute;
    \item Centerline-Modus \texttt{single}, Voxelgröße 0,1, B-Spline Grad 10,
      Eckensegmentierung deaktiviert.
\end{itemize}

  \subsection*{Ausführungsmodi}

  \begin{description}
    \item[Interaktiver Vollauf] führt durch Video-, COLMAP-, STS- und
      Meshroutenparameter sowie den optionalen GCP-Breakpoint.
    \item[Autopilot] setzt die dokumentierten Standardwerte und überspringt den
      manuellen GCP-Breakpoint, falls keine Matrix bereitsteht.
    \item[Replay] startet mit \texttt{--from} an einer Schrittgrenze und setzt
      konsistente vorhandene Artefakte voraus.
    \item[Matrixrunner] führt deklarierte Varianten seriell aus, archiviert sie
      und bereinigt danach den Live-Workspace.
  \end{description}

  Diese Modi dürfen in der Ergebnisinterpretation nicht gleichgesetzt werden.
  Insbesondere belegt eine erfolgreiche Matrix die Batchfähigkeit, aber nicht
  automatisch die vollständige Nutzerführung des Autopiloten.

  \subsection*{SuGaR-Fork: Versionierung und Detailnachweis}

  Der mask-aware Fork liegt unter \path{third_party/SuGaR}. Die für die
  Projektarbeit relevanten Änderungsschichten sind:

  \begin{itemize}
    \item \texttt{48bbfdd}: maskenbewusster Semantic-Mask-Provider sowie
      maskierte RGB-/DSSIM-Supervision in den Coarse-Trainern und im Refinement,
      maskierte DN-Consistency, konfigurierbare Maskenstufen/Dilatationen,
      semantisch begrenztes UV-Baking sowie erweiterte Trainings- und
      Meshparameter;
    \item \texttt{8f4f7a2}: Poisson-Schutz bei weniger als drei gültigen Punkten
      beziehungsweise ungültigen Normalen und Korrektur der
      Coarse-Iterationsprüfung für STS-Startzähler;
    \item \texttt{e254000}: Schalter für Gaussian-Rasterizer-Tiefe als
      kontrollierte Surface-Sampling-Diagnose;
    \item \texttt{ecda7ef}: zusätzliche Surface-Sampling-Optionen und
      explizite Background-Mesh-Steuerung;
    \item \texttt{eca4ea1} sowie der lokale Arbeitsbaum: normalisierte feste
      Eval-Splits, explizite Kardinalitätsprüfung und verständliche Fehler
      für nicht gematchte Kameras.
  \end{itemize}

  Die letzte Gruppe ist zum Zeitpunkt der Erstellung dieser Fassung noch nicht
  vollständig in den Parent-Gitlink übernommen. Vor der Abgabe müssen der
  Parent-Gitlink, der ausgecheckte Fork-Commit, der lokale Diff und
  \texttt{SUGAR\_REF} im finalen Manifest festgehalten werden.

  \subsection*{Bild- und Maskenablagen}

  Die Bildablage verwendet absichtlich mehrere Schnittstellen. SAM schreibt
  Arbeitsframes nach \path{data/02_frames}; COLMAP erzeugt bei einem
  verzeichnungsbehafteten Modell mit \texttt{image\_undistorter} eine physische
  ideale Kopie unter \path{data/04_sfm/undistorted/images}. STS erhält über
  \path{data/05_3dgs/images} nur einen Symlink auf diese ideale Szene. Die
  \path{data/03_masks/_attempts}-Ordner enthalten Prompt-Fallbacks; die flachen
  Masken sowie die drei hierarchischen Stufen sind Auswahl-, STS- und
  SuGaR-Vertragsdateien. Matrixarchive kopieren die relevanten Artefakte, damit
  ein Lauf unabhängig vom später bereinigten Live-Workspace replaybar bleibt.

  Die Reduktion dieser Ablagen ist eine Speicher-/I/O-Optimierung und kein
  elementarer Funktionsschritt. Elementar ist die inhaltliche Konsistenz von
  Frame, Kamera, idealem Bild und gewarpter Maske. Entfernt werden können nach
  erfolgreicher Auswahl vor allem nicht mehr benötigte Promptversuche oder
  regenerierbare Kompatibilitätskopien; die ideale Bilddomäne und der
  STS-Imagevertrag dürfen nicht ersatzlos gelöscht werden.

Die vollständige Laufhistorie mit Parametern, Status und Logs wird nicht in die
PA kopiert. Sie bleibt in den Matrixarchiven unter
\path{data/10_runs/}. Die zwölf abgeschlossenen SuGaR-Coarse-Läufe liegen
unter \path{data/10_runs/matrix_sugar_followup_12/}. Die Metriken der beiden
Modellstufen werden getrennt ausgewertet und über die gemeinsame CSV-Quelle in
\path{docs/grafiken/} nachvollziehbar. Vor einer finalen Abgabe sollten
Repository-Commit, SuGaR-Submodul-Commit, Docker-Image-Version und Hardware
zusätzlich in einem finalen Laufmanifest festgehalten werden.
